\section*{Hoofdstuk \ref{ch:g10:cells-common-unit} --- Cellen: de gemeenschappelijke eenheid van het leven}

\begin{solution}{exo:g10:cells-common-unit:1}
Elk levend wezen bestaat uit één of meer \omterm{def:g10:cells-common-unit:cell}{cellen}; de \omterm{def:g10:cells-common-unit:cell}{cel} is de
grondeenheid van bouw en werking; elke \omterm{def:g10:cells-common-unit:cell}{cel} komt door deling voort uit
een reeds bestaande \omterm{def:g10:cells-common-unit:cell}{cel}.
\end{solution}

\begin{solution}{exo:g10:cells-common-unit:2}
\qty{50}{\micro m}; \qty{2.5}{\micro m}; \qty{0.007}{mm}.
\end{solution}

\begin{solution}{exo:g10:cells-common-unit:3}
$36/600 = \qty{0.06}{mm} = \qty{60}{\micro m}$.
\end{solution}

\begin{solution}{exo:g10:cells-common-unit:4}
Alleen bij planten: \omterm{def:g10:cells-common-unit:organelle}{celwand}, \omterm{def:g10:cells-common-unit:organelle}{bladgroenkorrels}, grote centrale \omterm{def:g10:cells-common-unit:organelle}{vacuole}.
In elke \omterm{def:g10:cells-common-unit:cell}{cel}: een \omterm{def:g10:cells-common-unit:cell}{plasmamembraan} (en ook \omterm{def:g10:cells-common-unit:cell}{cytoplasma} en \omterm{def:g10:cells-common-unit:organelle}{ribosomen}).
\end{solution}

\begin{solution}{exo:g10:cells-common-unit:5}
Een \omterm{def:g10:cells-common-unit:prokaryote}{eukaryote cel} houdt haar DNA in een \omterm{def:g10:cells-common-unit:organelle}{celkern} en heeft door membranen
omsloten \omterm{def:g10:cells-common-unit:organelle}{organellen} (een wangcel, een gistcel); een \omterm{def:g10:cells-common-unit:prokaryote}{prokaryote cel} heeft
geen van beide, haar DNA ligt vrij in het \omterm{def:g10:cells-common-unit:cell}{cytoplasma}, en zij meet
ongeveer \qty{1}{\micro m} (een bacterie zoals \emph{E. coli}).
\end{solution}

\begin{solution}{exo:g10:cells-common-unit:6}
Werkelijke lengte $6 \times 5/15 = \qty{2}{\micro m}$. \omterm{def:g10:cells-common-unit:magnification}{Vergroting}: de
balk meet \qty{15}{mm} voor \qty{5}{\micro m} $= \qty{0.005}{mm}$, dus
$15/0.005 = 3000$: $\times 3000$.
\end{solution}

\begin{solution}{exo:g10:cells-common-unit:7}
Een \omterm{def:g10:cells-common-unit:organelle}{ribosoom} (\qty{25}{nm}) is tien keer kleiner dan het \omterm{def:g10:cells-common-unit:magnification}{oplossend
vermogen} van de lichtmicroscoop (\qty{0.2}{\micro m} $= \qty{200}{nm}$): twee
naburige \omterm{def:g10:cells-common-unit:organelle}{ribosomen}, of een \omterm{def:g10:cells-common-unit:organelle}{ribosoom} en zijn omgeving, zijn niet uit
elkaar te houden. Meer \omterm{def:g10:cells-common-unit:magnification}{vergroting} vergroot alleen de onscherpte.
\end{solution}

\begin{solution}{exo:g10:cells-common-unit:8}
De korrels zijn zetmeel, in de \omterm{def:g10:cells-common-unit:cell}{cellen} opgeslagen als vaste korrels
(zetmeel lost niet op). De voorraden van een plant liggen binnen in haar
\omterm{def:g10:cells-common-unit:cell}{cellen}, niet in een of andere ruimte ertussen.
\end{solution}

\begin{solution}{exo:g10:cells-common-unit:9}
De uienrok groeit ondergronds, in het donker, en slaat voorraden op:
\omterm{def:g10:cells-common-unit:organelle}{bladgroenkorrels} zouden daar nutteloos zijn. Het mosblaadje leeft in het
licht en maakt zijn voedsel door fotosynthese, dus zitten zijn \omterm{def:g10:cells-common-unit:cell}{cellen}
er vol mee. Hetzelfde plan, een andere uitrusting.
\end{solution}

\begin{solution}{exo:g10:cells-common-unit:10}
Ribbe \qty{1}{\micro m}: oppervlak \qty{6}{\micro m^2}, volume
\qty{1}{\micro m^3}, verhouding 6 per micrometer. Ribbe
\qty{20}{\micro m}: oppervlak \qty{2400}{\micro m^2}, volume
\qty{8000}{\micro m^3}, verhouding 0.3. De bacterie heeft twintig keer
meer membraan per eenheid inhoud en leeft alleen van diffusie; de
grotere \omterm{def:g10:cells-common-unit:cell}{cel} heeft inwendig transport en \omterm{def:g10:cells-common-unit:organelle}{organellen} nodig.
\end{solution}

\begin{solution}{exo:g10:cells-common-unit:11}
\qty{2}{h} zijn 6 delingen: $2^6 = 64$ \omterm{def:g10:cells-common-unit:cell}{cellen}. \qty{4}{h} zijn er 12:
$2^{12} = 4096$.
\end{solution}

\begin{solution}{exo:g10:cells-common-unit:12}
Zij heeft een membraan, een \omterm{def:g10:cells-common-unit:cell}{cytoplasma} en neemt stoffen op en geeft ze
af, dus voldoet zij aan het grootste deel van de definitie, maar zij kan
zich niet delen en haar \omterm{def:g10:chemistry-of-life:families}{eiwitten} niet vernieuwen: zij is een \omterm{def:g10:cells-common-unit:cell}{cel} die een
deel van het programma heeft opgegeven in ruil voor ruimte voor
hemoglobine, en zij wordt gemaakt en vervangen door volledige \omterm{def:g10:cells-common-unit:cell}{cellen}. De
definitie beschrijft het algemene geval; de rode bloedcel is daar een
gespecialiseerde eindvorm van.
\end{solution}

\begin{solution}{exo:g10:cells-common-unit:13}
Kurk is dood weefsel: Hooke zag de cellulosewanden van \omterm{def:g10:cells-common-unit:cell}{cellen} waarvan de
levende inhoud verdwenen was --- lege doosjes. Een \omterm{def:g10:cells-common-unit:organelle}{celkern} bestaat
alleen in het \omterm{def:g10:cells-common-unit:cell}{cytoplasma} van een levende \omterm{def:g10:cells-common-unit:cell}{cel}, en hoe dan ook zouden het
\omterm{def:g10:cells-common-unit:magnification}{oplossend vermogen} van zijn microscoop en het ontbreken van kleurstoffen
haar verborgen hebben.
\end{solution}

\begin{solution}{exo:g10:cells-common-unit:14}
Een virus heeft geen \omterm{def:g10:cells-common-unit:cell}{cytoplasma}, geen membraan van eigen makelij, geen
stofwisseling, en het kan zich niet delen: het voldoet niet aan de
definitie van een \omterm{def:g10:cells-common-unit:cell}{cel}. Volgens de celtheorie leeft het niet op eigen
kracht; het is een deeltje dat de machinerie van een \omterm{def:g10:cells-common-unit:cell}{cel} leent om
kopieën van zichzelf te maken. Of je het "levend" wilt noemen, is een
kwestie van afspraak; dat het geen \omterm{def:g10:cells-common-unit:cell}{cel} is, niet.
\end{solution}

\begin{solution}{exo:g10:cells-common-unit:15}
Volume $\pi \times (25\,\unit{\micro m})^2 \times 30000\,\unit{\micro m}
\approx \num{5.9e7}\,\unit{\micro m^3}$, ongeveer 7000 keer de
\qty{8000}{\micro m^3} van de kubus. Doordat zij een dunne cilinder is,
ligt geen enkel punt van de vezel verder dan \qty{25}{\micro m} van het
membraan, zodat de uitwisseling snel blijft; en elke kern bestuurt
alleen zijn eigen stuk van de vezel, zodat geen kern een groter volume
hoeft te bedienen dan dat van een gewone \omterm{def:g10:cells-common-unit:cell}{cel}.
\end{solution}

\begin{solution}{pb:g10:cells-common-unit:1}
\textbf{1.} $1.8/9 = \qty{0.2}{mm} = \qty{200}{\micro m}$.

\textbf{2.} $1.8/4 = \qty{0.45}{mm}$; er passen ongeveer twee \omterm{def:g10:cells-common-unit:cell}{cellen} over.

\textbf{3.} $28/0.2 = 140$: de tekening is $\times 140$.

\textbf{4.} $4/140 \approx \qty{0.029}{mm}$, ongeveer \qty{30}{\micro m}.

\textbf{5.} \omterm{def:g10:cells-common-unit:organelle}{Celwand}, \omterm{def:g10:cells-common-unit:cell}{plasmamembraan} (tegen de wand aan gedrukt),
\omterm{def:g10:cells-common-unit:cell}{cytoplasma}, \omterm{def:g10:cells-common-unit:organelle}{celkern}, \omterm{def:g10:cells-common-unit:organelle}{vacuole}. Wel aanwezig maar onzichtbaar bij
$\times 400$: de \omterm{def:g10:cells-common-unit:organelle}{ribosomen} (en de \omterm{def:g10:cells-common-unit:organelle}{mitochondriën} liggen op de grens).

\textbf{6.} Zij komen van het oppervlak van een bekleding die enkele
\omterm{def:g10:cells-common-unit:cell}{cellen} dik is en voortdurend wordt geschuurd: de buitenste \omterm{def:g10:cells-common-unit:cell}{cellen} zijn
afgeplat en laten los, en daarom vangt een wattenstaafje ze op.

\textbf{7.} \qty{0.24}{mm} ligt net boven de grens van \qty{0.1}{mm} van
het blote oog: een nauwelijks zichtbaar stipje. Bij $\times 100$ beslaat
het $0.24/1.8 \approx 13\%$ van het veld, bij $\times 400$ meer dan de
helft; geen enkel beschikbaar objectief geeft precies een kwart.

\textbf{8.} Bacteriën. De elektronenmicroscoop zou hun wand, membraan,
\omterm{def:g10:cells-common-unit:organelle}{ribosomen} en DNA-gebied tonen, en het ontbreken van een \omterm{def:g10:cells-common-unit:organelle}{celkern}.

\textbf{9.} De wangcel doet één taak (een oppervlak beschermen) in een
lichaam waar andere \omterm{def:g10:cells-common-unit:cell}{cellen} haar voeden, van zuurstof voorzien en
aansturen; het pantoffeldiertje heeft niemand anders en moet als één
\omterm{def:g10:cells-common-unit:cell}{cel} eten, bewegen, waarnemen en zich delen.

\textbf{10.} De buitenkant is nu zouter dan het \omterm{def:g10:cells-common-unit:cell}{cytoplasma}, dus verlaat
water de \omterm{def:g10:cells-common-unit:cell}{cel} door het membraan; de \omterm{def:g10:cells-common-unit:cell}{cel} verliest volume en schrompelt
ineen.

\textbf{11.} $20^3 = \qty{8000}{\micro m^3} = \num{8e-6}\,\unit{mm^3}$.

\textbf{12.} $6 \times 20^2 = \qty{2400}{\micro m^2}$; verhouding
$2400/8000 = 0.3$ per micrometer.

\textbf{13.} Oppervlak \qty{6}{\micro m^2}, volume \qty{1}{\micro m^3},
verhouding 6: twintig keer groter.

\textbf{14.} $\qty{8000}{\micro m^3} = \num{8e-9}\,\unit{cm^3}$, dus
ongeveer $\num{8e-9}\,\unit{g} = \qty{8}{ng}$.

\textbf{15.} $\pi \times 3.5^2 \times 2 \approx \qty{77}{\micro m^3}$,
ongeveer een honderdste van de kubus.

\textbf{16.} $\qty{60}{L} = \num{6e16}\,\unit{\micro m^3}$; gedeeld door
8000: ongeveer $\num{7.5e12}$ \omterm{def:g10:cells-common-unit:cell}{cellen}.

\textbf{17.} Twee derde daarvan: ongeveer $\num{5e12}$.

\textbf{18.} $\num{2.5e13} \times 77 \approx \num{1.9e15}\,\unit{\micro
m^3} \approx \qty{1.9}{L}$. Bloedcellen: $0.45 \times 4.5 \approx
\qty{2.0}{L}$ --- dat klopt.

\textbf{19.} Ongeveer $\num{5e12}$ grote \omterm{def:g10:cells-common-unit:cell}{cellen} plus $\num{2.5e13}$ rode
\omterm{def:g10:cells-common-unit:cell}{cellen}: zo'n $\num{3e13}$, enkele tientallen biljoenen. Celgroottes
verschillen een factor honderd en de verhouding tussen de typen
wisselt, dus elke telling is een schatting op een factor twee of drie na.

\textbf{20.} Ruwweg dertig biljoen ($\num{3e13}$) \omterm{def:g10:cells-common-unit:cell}{cellen}, alle ontstaan
door opeenvolgende celdelingen uit één beginnende \omterm{def:g10:cells-common-unit:cell}{cel}, de bevruchte
eicel.
\end{solution}
